% Teil II: durchgerechneter Weg, Teile a--d.

\subsection*{Gegeben und gesucht}
$m=\qty{0.1}{kg}$, $R_1=\qty{30}{cm}=0{,}30\,\mathrm m$,
$\omega_1=\qty{9}{s^{-1}}$, $R_2=\qty{15}{cm}=0{,}15\,\mathrm m$; global
f{\"u}r $r_1=R_1$, $r_2=R_2$. Gesucht: (a) Erhaltung und $\omega_2$,
(b) Fadenspannung bei $R_1$ und $R_2$, (c) Arbeit der ziehenden Kraft,
(d) Aenderung der kinetischen Energie und Vergleich.

\subsection*{Ansatz}
Drehimpulserhaltung $mr^2\omega=\text{const.}$; Fadenspannung =
Zentripetalkraft; Arbeit = Energiebilanz $\Delta E_{\rm kin}$.

\subsection*{Rechnung}
\textbf{a) Erhaltung und $\omega_2$.} Die Zugkraft wirkt entlang des
Fadens, also radial durch die Mitte -- sie uebt kein Drehmoment
bezueglich der Achse aus, deshalb bleibt $L=mr^2\omega$ erhalten.
\begin{equation}
  \omega_2=\omega_1\left(\frac{R_1}{R_2}\right)^2
  =\qty{9}{s^{-1}}\left(\frac{0{,}30}{0{,}15}\right)^2
  =\qty{9}{s^{-1}}\cdot4
  =\boxed{\;\qty{36}{s^{-1}}\;}.
\end{equation}

\textbf{b) Fadenspannung.}
\begin{equation}
  F=m\,r\,\omega^2: \quad
  F_1=mR_1\omega_1^2=0{,}1\cdot0{,}30\cdot81\approx\boxed{\;\qty{2.43}{N}\;},
\end{equation}
\begin{equation}
  F_2=mR_2\omega_2^2=0{,}1\cdot0{,}15\cdot1296
  \approx\boxed{\;\qty{19.4}{N}\;}
\end{equation}

\textbf{c) Arbeit der ziehenden Kraft.} $\Delta E_{\rm kin}$, mit
$E_{\rm kin}=\tfrac12 m r^2\omega^2=\tfrac{L^2}{2mr^2}$:
\begin{align}
  E_1&=\tfrac12 m R_1^2\omega_1^2
      =0{,}5\cdot0{,}1\cdot0{,}09\cdot81\approx\qty{0.365}{J},\\
  E_2&=\tfrac12 m R_2^2\omega_2^2
      =0{,}5\cdot0{,}1\cdot0{,}0225\cdot1296\approx\qty{1.46}{J},\\
  \Delta E&=E_2-E_1\approx\qty{1.46}{J}-\qty{0.365}{J}
      =\boxed{\;\qty{1.09}{J}\;}.
\end{align}

\textbf{d) Vergleich.} Die kinetische Energie aendert sich genau um den
Betrag der Arbeit der Zugkraft:
\begin{equation}
  \boxed{\;W=\Delta E_{\rm kin}=\qty{1.09}{J}\;}.
\end{equation}
Die Energie steigt an (der Faden zieht nach innen, die Zentripetalkraft
hat eine radiale Komponente; die Zugkraft verrichtet Arbeit).

\begin{ergebnis}
$\omega_2=\qty{36}{s^{-1}}$; $F_1\approx\qty{2.43}{N}$,
$F_2\approx\qty{19.4}{N}$; $W=\Delta E_{\rm kin}\approx\qty{1.09}{J}$;
die kinetische Energie waechst um diesen Betrag.
\end{ergebnis}

\begin{kontrolle}
\textbf{Einheiten:} $\omega$ [s$^{-1}$]; $F=mr\omega^2$: kg$\cdot$m$\cdot$s$^{-2}$
[$\rightarrow$N]. $E=\tfrac12mr^2\omega^2$: kg$\cdot$m$^2\cdot$s$^{-2}$=J. \checkmark

\textbf{Querprobe:} Radius halbiert sich, also $\omega_2=4\omega_1$;
$E\propto r^2\omega^2\propto r^2(1/r^4)=1/r^2$: bei halbem Radius vierfache
Energie. Von $E_1=\qty{0.365}{J}$ zu $E_2\approx4\cdot0{,}365=1{,}46\,\mathrm J$ --
stimmt. \checkmark

\textbf{Spannung:} $F_2\approx8F_1$ (da $\omega_2^2R_2/(\omega_1^2R_1)
=16\cdot\frac12=8$) -- Anstieg deutlich. \checkmark
\end{kontrolle}

\begin{merksatz}
Zieht man eine Masse auf der Kreisbahn nach innen, bleibt der Drehimpuls
erhalten ($\omega\propto1/r^2$); die Arbeit der Zugkraft gleicht dem
Energiezuwachs $E_{\rm kin}=L^2/(2mr^2)$ und nimmt mit kleinerem Radius zu.
\end{merksatz}
